\clearpage
\thispagestyle{fancy}
\begin{center}
    {\Large \textbf{Arabic Summary}} 
\end{center}

% Move the spacing command here, outside the Arabic block
\setlength{\parskip}{1em}

\begin{arabtext}
\setcode{utf8}
تقدم هذه الرسالة نظاماً تجميعياً روبوتياً تعاونياً يدمج بين ذراع صناعية من طراز (\L{ABB}) وذراع روبوتية من طراز (\L{AR4}) لتنفيذ عملية تجميع الطوب بكفاءة، وذلك من خلال تنسيق عمليات الإدراك، والقبض، والتسليم، والوضع. يجمع النظام المقترح بين الاستيعاب الشامل للمشهد عبر كاميرا (\L{Intel RealSense}) وبين الاستشعار المحلي (\L{eye-in-hand}) باستخدام كاميرا (\L{ESP32-CAM}) مثبتة على ذراع (\L{AR4})، مما يتيح تحديد المواقع التقريبي والمحاذاة الدقيقة أثناء القبض عبر نظام التوجيه البصري القائم على الصور (\L{IBVS}). وتقوم "آلة حالة" (\L{State Machine}) مهيكلة بإدارة تخصيص المهام وتنفيذها من البداية إلى النهاية، حيث تختار بين التجميع المباشر أو التسليم المتبادل بناءً على منطق تخصيص الطوب والروبوت المستمد من مدخلات واجهة المستخدم الرسومية (\L{GUI})، وقاعدة البيانات، والإدراك الشامل للمشهد.

لتحقيق تفاعل دقيق في مساحة عمل مشتركة، تم التعامل مع المعايرة كمتطلب أساسي. أُجريت معايرة ذراع (\L{AR4}) عبر جمع البيانات عند مواقع (\L{X--Y}) متعددة وعلى مستويين من الارتفاع (\L{0 mm} و\L{50 mm})، مع مطابقة دوال الربط باستخدام طرق متعددة (دالة القاعدة الشعاعية (\L{RBF})، الخطية، والمتعددة الحدود). وقد أظهر الانحدار متعدد الحدود أداءً فائقاً وتم التحقق من صحته باستخدام تقسيمات بيانات التدريب والاختبار. وبالمثل، تم تطبيق معايرة الكاميرا الشاملة بناءً على الانحدار متعدد الحدود باستخدام عينات عبر إحداثيات \L{X, Y, Z} متنوعة لتمكين التقدير الموثوق لوضعية الطوب في العالم الحقيقي عبر الطاولة بأكملها. كما تم تنفيذ الإدراك البصري من خلال خوارزمية (\L{YOLOv8}) لتقسيم الصور واكتشاف الطوب والمعالم، بينما تمت معالجة اختيار القبضة باستخدام "محول" (\L{Transformer}) قائم على الترميز وفك التشفير، تم تدريبه على حالات القبض الناجحة (\L{True Grasps}) وحالات الفشل/التجنب (\L{False Grasps})، مما أتاح استخدام التعزيز السلبي لتقليل تنبؤات القبض غير المستقرة. وأثناء التحكم الدقيق، تم استخدام "التماثل البياني" (\L{Homography}) لربط نقاط المعالم بين الرؤية الشاملة لكاميرا (\L{RealSense}) والرؤية المحلية لكاميرا (\L{ESP32-CAM})، مما بسّط عملية التتبع وحسّن من موثوقية الاقتراب.
\end{arabtext}